% ch34-fp8-quant-gemm.tex — FP8 GEMM（一）：e4m3 格式、per-row/per-block 量化与前处理 Kernel（RFC-N）
% 源码：kernels/interview/fp8_gemm.cuh（Phase 9.1/9.2，L45-284）
\chapter{FP8 GEMM（一）：e4m3 格式、per-row/per-block 量化与前处理 Kernel}\label{ch:34}

\section{本章导读}

第~\ref{ch:27}--\ref{ch:33}~章（第六部分）回答了「量化 attention 怎么做」；本部分
把同一套 e4m3 数学搬到\textbf{矩阵乘}上，拆成两章：本章写\textbf{量化前处理
kernel}——把 BF16 的 $A,B$ 变成 e4m3 的 $A_8,B_8^{\top}$ 外加两组 scale；
第~\ref{ch:35}~章写\textbf{主 GEMM kernel 与在线反量化 epilogue}，并把全链路封装成
一个 C++ API 与 cuBLAS BF16 GEMM 对精度和性能。两章合起来构成的完整调用形态是：

\begin{equation}
  \underbrace{A,B}_{\text{BF16 输入}}
  \;\xrightarrow[\text{本章：3 个量化 kernel}]{\text{amax 统计 + 量化 + 转置}}
  \underbrace{A_8,\ B_8^{\top},\ \delta_a,\ \delta_b}_{\text{FP8 workspace}}
  \;\xrightarrow[\text{第\ref{ch:35}章：CuTe 主 kernel}]{\text{TMA + mma.sync + epilogue}}
  \underbrace{C}_{\text{BF16 输出}}
\end{equation}

与 attention 量化相比，GEMM 量化有一处本质的简化：\textbf{scale 折叠可以完全
离开主循环}。attention 的量化输出要参与 softmax，scale 必须在 kernel 内部按行
修正（第~\ref{ch:29}~章的 $s_{dequant}=q_s k_s$ 折叠也要挂在每个 epilogue 上）；
而 GEMM 是双线性运算，两组 scale 可以\textbf{推迟到 epilogue 一次乘回}，主循环
里跑的是纯整数域（e4m3 编码值）的 MMA。这个差异是 34.4~节的核心恒等式，也是
第~\ref{ch:35}~章 epilogue 设计的全部依据。

前处理链本身是三个「读一遍、写一遍」的流式 kernel（带宽型，不是算力型），工程
重点全在\textbf{粒度}与\textbf{布局}两件事上：
\begin{enumerate}[itemsep=0pt,topsep=2pt]
  \item \textbf{粒度}：A 支持 per-row（每行一个 scale，精度最优）与 per-block
        （每 128 行一个，ffpa Q/K 同构）；B 支持per-col 与 per-block。四种组合
        全部实现（34.4~节），精度开销各不同；
  \item \textbf{布局}：$A_8$ 保形行主序；$B$ 必须转置量化成 $B_8^{\top}[N,K]$
        行主序——主 kernel 的 TMA 与 MMA 都按 TN 布局消费，转置发生在量化
        kernel 内部（smem staging，34.5.3~节），对调用方完全透明。
\end{enumerate}

\textbf{前置章节}：第~\ref{ch:27}~章（e4m3 编码与量化算子，必须）、
第~\ref{ch:02}~章（向量化 IO 与 warp 归约）、第~\ref{ch:03}~章（block 归约）。
\textbf{源码区间}：\texttt{fp8\_gemm.cuh} L45--284（Phase 9.1/9.2），完整行号以
附录~\ref{app:D}~的源码索引为准。
\textbf{验证方式}：\texttt{notes\_v2\_sm120a.bin --fp8-gemm}（CPU fp64 参考全
组合 + 尾部 shape），bench 数据在 34.6~节。

\section{问题与动机}

\subsection{为什么 GEMM 量化比 attention 量化「容易」}

attention 的三量化点里，$P$（softmax 输出）必须在线量化，$Q,K,V$ 的 scale 又会
经 softmax 混入归一化分母——所以第~\ref{ch:29}~章要在 kernel 里维护行最大值
修正与 $s_{dequant}$ 折叠。GEMM 没有任何非线性：$C=AB^{\top}$ 是双线性的，
量化的 scale 是\textbf{纯标量因子}，提到括号外就行。于是：

\begin{itemize}[itemsep=1pt,topsep=2pt]
  \item 主 kernel 的 MMA 循环里\textbf{一个 scale 都不用碰}，跑的就是 e4m3 编码
        值的整数语义乘累加（f32 累加器）；
  \item 反量化只在 epilogue 出现一次：$C = (A_8 B_8^{\top})\odot(\delta_a\delta_b)$，
        $\odot$ 在 per-block 粒度下退化为块级（128 行/列）标量乘（第~\ref{ch:35}~章）；
  \item 前处理全部离线于主 kernel：量化 kernel 写完 workspace，主 kernel 才进场。
\end{itemize}

教学上这是一条「从第~\ref{ch:27}~章的量化数学到工程 kernel」的最短路径：数学
恒等式一章讲透，kernel 工程两章讲透，不夹带 online softmax 这类复杂度。

\subsection{约束与 API 契约}

整套实现的形状约束由\textbf{TMA 对齐}与\textbf{量化向量化}共同决定，集中在 API
层检查（第~\ref{ch:35}~章 \texttt{fp8\_gemm\_bf16}）：

\begin{itemize}[itemsep=1pt,topsep=2pt]
  \item $K \bmod 16 = 0$：$A_8/B_8^{\top}$ 的行是 e4m3 字节流，TMA 装载要求
        16\,B 对齐；同时也是量化 kernel Vec8（8$\times$bf16$=$16\,B）读取的
        超集；
  \item $N \bmod 8 = 0$：输出 $C$ 的行是 bf16 字节流，TMA store 的 16\,B 对齐
        要求 $N\times 2$\,B 被 16 整除；
  \item $M$ 任意：装载方向的越界行由 TMA 自动零填充（对累加无贡献），写出
        方向越界行由 TMA 自动裁剪，epilogue 的 scale 查表按 $m<M$ guard；
  \item $K$ 非 128 倍数的尾 tile 同样零填充，主循环\textbf{没有尾部分支}——
        这比第~\ref{ch:28}~章 ffpa 的 VT 量化省掉了 pad 列清零，因为那里用的是
        cp.async 装载、这里全程 TMA（OOB 语义内建）。
\end{itemize}

\section{e4m3 回顾：GEMM 视角}\label{sec:34-e4m3}

第~\ref{ch:27}~章已逐位讲过 e4m3 的编码；这里只保留 GEMM 工程要用的三个事实，
并用一张位布局图统一符号。e4m3 是 1 符号 + 4 指数 + 3 尾数的 8 位浮点：指数
偏置 7、无无穷大（全 1 指数域 + 非零尾数是 NaN），最大有限值 $R=448$：

\begin{figure}[H]
\centering
\begin{tikzpicture}[
  font=\footnotesize,
  bit/.style={draw, thick, minimum width=0.92cm, minimum height=0.92cm, inner sep=0pt},
  lbl/.style={font=\scriptsize, anchor=south},
]
% bit boxes: s e e e e m m m
\foreach \i/\c in {0/s,1/e,2/e,3/e,4/e,5/m,6/m,7/m} {
  \node[bit] at (\i*0.95, 0) {\c};
}
\node[lbl] at (0*0.95, 0.62) {符号 $s$};
\node[lbl] at (2.5*0.95, 0.62) {指数 $E$（4 bit，偏置 7）};
\node[lbl] at (6*0.95, 0.62) {尾数 $M$（3 bit）};
\node[font=\scriptsize, anchor=north west, align=left] at (-0.55,-0.62)
  {数值 $=(-1)^s\times 2^{E-7}\times (1{+}M/8)$\\
   规格化域 $2^{-6}\sim 448$，相对步长 $2^{-3}$（半步舍入误差 $\le 2^{-4}$）\\
   饱和上界 $R=448$：\texttt{cvt.rn.satfinite} 越界饱和不产生 NaN};
\end{tikzpicture}
\caption{e4m3 位布局与 GEMM 工程三要素：相对步长固定 $2^{-3}$（尾数 3 位），
决定量化噪声的下限；满量程 448 决定 scale 公式 $\delta=\mathrm{amax}/448$；
\texttt{SATFINITE} 语义决定量化 kernel 可以放心写 $x\cdot\delta^{-1}$ 而不用
先 clamp——硬件饱和替我们做了。}
\label{fig:34-e4m3}
\end{figure}

三个事实：
\begin{enumerate}[itemsep=0pt,topsep=2pt]
  \item \textbf{相对步长 $2^{-3}$}：规格化域内相邻可表示值的间距是值本身的
        $1/8$，round-to-nearest 的单元素相对误差上界 $2^{-4}=6.25\%$。这决定了
        34.4~节的统计误差模型；
  \item \textbf{满量程 $R=448$}：对称量化的 scale 定义为
        $\delta = \mathrm{amax}/448$，让块内最大绝对值恰好映到编码域边缘，
        步长利用率最大；
  \item \textbf{饱和语义}：\texttt{\_\_nv\_cvt\_float2\_to\_fp8x2} 的
        \texttt{\_\_NV\_SATFINITE} 模式把越界输入饱和到 $\pm448$ 而不是
        NaN——量化 kernel 不需要软件 clamp（34.5.4~节）。
\end{enumerate}

\section{量化数学与误差分析}\label{sec:34-math}

\subsection{量化算子与 scale 折叠恒等式}

对称量化算子（与第~\ref{ch:27}~章式~(27.x) 同一形式，重写于此保持自包含）：

\begin{equation}
  \hat{x} = Q(x;\delta) = \mathrm{sat}_{[-448,448]}\!\big(\mathrm{rn}(x/\delta)\big),
  \qquad \delta = \frac{\mathrm{amax}(\mathcal{B})}{448},
  \label{eq:34-quant-op}
\end{equation}
其中 $\mathcal{B}$ 是 scale 的统计块（一行 / 一列 / 一个 $128$ 行块，由粒度
决定），$\mathrm{rn}$ 是 round-to-nearest-even。反量化 $x \approx \hat{x}\cdot\delta$。

GEMM 的核心恒等式由此立得。设 $A \approx \delta_a \hat{A}$，
$B \approx \delta_b \hat{B}$（$\delta_a,\delta_b$ 为逐行/逐列/逐块的 scale
张量， broadcasting 到元素级），则

\begin{equation}
  C = A B^{\top} = (\delta_a \odot \hat{A})(\delta_b \odot \hat{B})^{\top}
    \;\neq\; \text{别的形态，}
  \label{eq:34-fold}
\end{equation}

精确展开单元素：当 $\delta_a$ 沿行（$m$ 方向）常数、$\delta_b$ 沿列（$n$ 方向）
常数时——per-row $\times$ per-col 正是如此——
\begin{equation}
  C_{m,n} = \sum_{k} \delta_{a,m}\,\hat{A}_{m,k}\,\hat{B}_{n,k}\,\delta_{b,n}
          = \delta_{a,m}\,\delta_{b,n} \sum_k \hat{A}_{m,k}\hat{B}_{n,k}
          = \delta_{a,m}\,\delta_{b,n}\,(\hat{A}\hat{B}^{\top})_{m,n}.
  \label{eq:34-fold-rc}
\end{equation}
即\textbf{scale 完全出括号}：整数域（e4m3 编码值）MMA 之后，epilogue 按
$(m,n)$ 查 $\delta_{a,m}\delta_{b,n}$ 一次乘回。per-block 粒度时块宽固定 128
（与 tile 尺寸无关），tile 内 $\delta_a,\delta_b$ 是\textbf{分段常数}：$k_{BM}{=}128$
的 A 侧一个 tile 恰一块，$k_{BN}{=}256$ 的 B 侧一个 tile 横跨两块，故查表索引
取全局的 $m/128$、$n/128$，式~\ref{eq:34-fold-rc} 退化为\textbf{块级标量乘}。

对比第~\ref{ch:29}~章 persist-D 的 $s_{dequant}=q_s k_s$：那里 Q/K 的 scale 同样
出括号，但 V 的 scale 乘在 $P$ 上、而 $P$ 又来自 softmax，折叠点必须跟着
在线重缩放走；GEMM 版没有这层耦合，恒等式干净得多。这也是本部分先于
attention 部分教学的价值：先看清「纯双线性下量化能折叠到什么程度」。

\subsection{四种粒度组合}

实现里 \texttt{Fp8GemmScaleMode} 枚举出四种组合，全部走同一对恒等式，差异只在
scale 张量的形状与 epilogue 的查表方式：

\begin{figure}[H]
\centering
\begin{tikzpicture}[
  font=\footnotesize,
  cell/.style={draw=tkPurple!60, minimum width=0.30cm, minimum height=0.30cm, inner sep=0pt},
  grpA/.style={draw=tkBlue, very thick},
  grpB/.style={draw=tkGreen, very thick},
  note/.style={font=\scriptsize, align=left},
]
% --- A: per-row ---
\begin{scope}[shift={(0,0)}]
  \foreach \r in {0,...,5}
    \foreach \c in {0,...,7}
      \node[cell] at (\c*0.30, -\r*0.30) {};
  \draw[grpA] (-0.21, 0.21) rectangle (2.31, -0.21);
  \draw[grpA] (-0.21, -0.99) rectangle (2.31, -1.41);
  \node[note, anchor=north] at (1.2, -1.85) {A per-row：$\delta_a[m]$，共 $M$ 个\\{\tiny 每行 amax，精度最优，scale 读写 $O(M)$}};
\end{scope}
% --- A: per-block ---
\begin{scope}[shift={(4.1,0)}]
  \foreach \r in {0,...,5}
    \foreach \c in {0,...,7}
      \node[cell] at (\c*0.30, -\r*0.30) {};
  \draw[grpA, rounded corners=1pt] (-0.21, 0.21) rectangle (2.31, -1.71);
  \node[note, anchor=north] at (1.2, -1.85) {A per-block：$\delta_a[m/128]$，$M/128$ 个\\{\tiny 整块单一 amax（ffpa Q/K 同构），outlier 摊最粗}};
\end{scope}
% --- B: per-col(列方向框) ---
\begin{scope}[shift={(0,-4.3)}]
  \foreach \r in {0,...,5}
    \foreach \c in {0,...,7}
      \node[cell] at (\c*0.30, -\r*0.30) {};
  \draw[grpB] (-0.21, 0.21) rectangle (0.21, -1.71);
  \draw[grpB] (0.99, 0.21) rectangle (1.41, -1.71);
  \node[note, anchor=north] at (1.2, -1.85) {B per-col：$\delta_b[n]$，共 $N$ 个\\{\tiny 每列 amax（实现转置为 B8T 按行量化）}};
\end{scope}
% --- B: per-block ---
\begin{scope}[shift={(4.1,-4.3)}]
  \foreach \r in {0,...,5}
    \foreach \c in {0,...,7}
      \node[cell] at (\c*0.30, -\r*0.30) {};
  \draw[grpB, rounded corners=1pt] (-0.21, 0.21) rectangle (2.31, -1.71);
  \node[note, anchor=north] at (1.2, -1.85) {B per-block：$\delta_b[n/128]$，$N/128$ 个\\{\tiny 128 列共用一个 amax}};
\end{scope}
\end{tikzpicture}
\caption{四种 scale 粒度组合（上排 A，下排 B）。蓝框=A 侧 scale 统计范围，
绿框=B 侧。粒度越粗，scale 存储/查表越省（per$\times$per-block 每 tile 只查
两个标量），但 outlier 的「毒性」越大——一个大值会把整块的步长撑爆，
其余小值全部挤进低位。四种组合在 API 里是一个枚举，主 kernel 用
\texttt{if constexpr} 编译期剔除，运行时零分支开销。}
\label{fig:34-granularity}
\end{figure}

为什么 GEMM 里 per-row$\times$per-col 是「精度最优」：scale 的统计块越小，块内
分布越均匀、步长越贴身。极端的反例是全矩阵单 scale（per-tensor）：一个 outlier
可以让普通元素的相对量化误差放大若干个数量级。Per-block 是 attention 领域
（SageAttention/ffpa）的主流选择，因为它与 128 的 tile 尺寸天然对齐、kernel 内
scale 折成单标量；对 GEMM 的权重矩阵（分布沿行/列明显不均），per-row/per-col
通常明显更准。34.6~节会给四种组合的同机实测。

\subsection{误差模型：能差到多少}\label{sec:34-err}

单元素量化噪声：设块内元素落在同一二进区间，$\mathrm{rn}$ 的误差均匀分布于
$[-\frac{u}{2}, \frac{u}{2}]$（$u=2^{-3}$ 为相对步长），故单元素相对误差
$\varepsilon$ 满足 $|\varepsilon|\le 2^{-4}$，方差 $\sigma_\varepsilon^2 = u^2/12$。

scale 折叠之后，误差只来自编码舍入一个源头。把单元素量化写成乘性形式
$\hat{x} = x(1+\varepsilon)$：相对舍入误差 $\varepsilon$ 在
$[-2^{-4}, 2^{-4}]$ 上近似均匀分布（相对步长 $u = 2^{-3}$，见上文），故
$\mathrm{E}[\varepsilon] \approx 0$，方差 $\sigma_\varepsilon^2 = u^2/12
\approx 1.30\times 10^{-3}$，即 $\sigma_\varepsilon \approx 3.6\%$。
对 $C_{m,n}=\sum_k a_k b_k$，两个量化点的误差一阶展开：
\begin{equation}\label{eq:34-err-expand}
  \tilde{C}_{m,n} - C_{m,n}
  = \sum_{k} a_k b_k\bigl(\varepsilon_{a,k} + \varepsilon_{b,k}\bigr)
  + \sum_k a_k b_k\,\varepsilon_{a,k}\varepsilon_{b,k}
  \;\approx\; \sum_{k} a_k b_k\bigl(\varepsilon_{a,k} + \varepsilon_{b,k}\bigr),
\end{equation}
第二项是 $O(\sigma_\varepsilon^2)$ 的高阶小量，比一阶项低一个
$\sigma_\varepsilon$ 因子，下文只保留一阶。关键区分两种数据形态：
\begin{itemize}[itemsep=1pt,topsep=2pt]
  \item \textbf{随机符号数据}（测试用的均匀随机矩阵）：误差项相互独立且
        随机变号，累加是随机游走。误差能量
        $\mathrm{Var}[\tilde{C}-C] = K\cdot\mathrm{E}[(ab)^2]\cdot
        2\sigma_\varepsilon^2$，而 $C$ 自身能量
        $\mathrm{Var}[C] = K\cdot\mathrm{E}[(ab)^2]$，相对误差
        $\sqrt{2}\,\sigma_\varepsilon \approx 5\%$（实测 3.6\%，
        34.6~节——舍入误差并非全程满幅均匀，实测略低于上界）。
        注意相对误差\textbf{与 $K$ 无关}：误差与信号都以 $\sqrt{K}$
        增长，抵消掉了；
  \item \textbf{同号数据}（全正矩阵、经 ReLU 的激活）：误差项同号不相消，
        累加是线性增长，最坏相对误差停在单元素上界 $2^{-4}$ 量级附近
        （不随时长收敛）。这就是量化 GEMM 在真实模型上仍需
        per-row/per-channel 粒度与离线校准（如 GPTQ/平滑旋转，
        第~\ref{ch:28}~章同源思想）的原因。
\end{itemize}
\paragraph{Max Err 是极值统计, 不是误差上界.} 单个 $C_{m,n}$ 接近零（正负
相消）时相对误差可以任意大，所以 bench 表里的 \texttt{Max Err} 只能取绝对
误差（vs cuBLAS BF16 参照）。但它也不是最坏元素误差：$MN$ 个近似独立的
误差样本取最大值，极值统计给出
\begin{equation}\label{eq:34-err-max}
  \mathrm{E}\bigl[\max_{m,n} |\tilde{C}-C|\bigr]
  \;\approx\; \sigma_{\text{err}}\sqrt{2\ln(MN)},
\end{equation}
$M{=}N{=}4096$ 时 $\sqrt{2\ln(MN)} \approx 5.8$。误差分布是有限支撑的和
分布，尾部比高斯瘦，实测极值落在 $4\sigma_{\text{err}}$ 一带——健康判据是
\textbf{Max Err 在 $4\sigma_{\text{err}}$ 量级}，而不是要求它趋零。绝对
口径的 Max Err 自然随 $\sqrt{K}$ 增长（$\sigma_{\text{err}} \propto
\sqrt{K}$），换算到相对口径才是与 $K$、与分布都无关的常数。

\paragraph{实测指纹.} 四组数据把上面每个环节都钉住（$\sigma_{\text{err}}
= \sqrt{K\,\mathrm{E}[(ab)^2]\,2\sigma_\varepsilon^2}$，均匀输入
$\mathrm{E}[(ab)^2] = 1/9$，randn 组 $\mathrm{E}[(ab)^2] = \sigma_x^4$，
$\sigma_x = 0.25/3$）：

\begin{center}
\footnotesize
\begin{tabular}{lccc}
\toprule
数据形态（vs cuBLAS BF16） & $\sigma_{\text{err}}$ 理论 & Max Err 实测 & 比值 \\
\midrule
均匀 $[-1,1]$，2048$^3$ & 0.77 & 3.00 & 3.9$\sigma$ \\
均匀 $[-1,1]$，4096$^3$ & 1.09 & 4.50 & 4.1$\sigma$ \\
均匀 $[-1,1]$，8192$^3$ & 1.54 & 6.25 & 4.1$\sigma$ \\
randn $(\pm 0.25)$，4096$^3$ & 0.0227 & 0.090 & 4.0$\sigma$ \\
\bottomrule
\end{tabular}
\end{center}

三组均匀 shape 的 Max Err/$\sigma_{\text{err}}$ 停在同一个 $4.1\sigma$——
样本量从 $4\text{M}$ 涨到 $67\text{M}$，比值不漂移，证明误差是零均值随机
游走而非系统性偏差。randn 组把输入压成 $\sigma_x = 0.25/3$ 的截断正态
（幅值缩 $s \approx 6.9$ 倍）：$\mathrm{E}[(ab)^2]$ 缩 $s^4$，
$\sigma_{\text{err}}$ 随之缩 $s^2 \approx 48$ 倍，$4.50/48 \approx 0.094$
对上实测 0.090，比值仍钉在 $4\sigma$。注意两组的相对误差
$\sigma_{\text{err}}/\sigma_C = \sqrt{2}\,\sigma_\varepsilon$ 相同：换分布
只缩放绝对口径，34.6~节实测的 rel-Frobenius 0.036 与分布无关。

\paragraph{与 fp8 attention 的误差为什么不可直接比.} 同是 fp8 对 bf16，
量化 attention 的绝对误差看起来小得多，差异全在算子形态与度量口径。其一，
attention 输出 $O = \mathrm{softmax}(QK^{\top})V$ 是 $V$ 行的\textbf{凸
组合}，元素被夹在 $[\min V, \max V]$ 之间，误差也被加权平均冲淡；GEMM 的
$C$ 是 $K$ 项\textbf{无归一化}累加，$\sigma_C = \sqrt{K\,\mathrm{E}[(ab)^2]}$
随 $\sqrt{K}$ 膨胀。其二，attention 惯用 PSNR/相对误差口径（几个百分点），
bench 表用绝对误差口径。同口径对比：本实现 rel-Frobenius 3.6\% 与 fp8
attention 对 bf16 的相对误差同量级，没有谁比谁误差大一说。

\subsection{amax 零块与数值安全}

$\mathrm{amax}=0$（全零行/块）时 $\delta=0$、$x/\delta$ 为 $0/0$。实现的约定：
\begin{equation}
  \delta = 0,\qquad \delta^{-1} \mathrel{\hat=} 0
  \;\Longrightarrow\; \hat{x} = \mathrm{rn}(x\cdot\delta^{-1}) = 0,
\end{equation}
即全零块量化为全零编码（合法 e4m3 值），反量化 $0\times 0 = 0$ 精确成立。
这个 \texttt{amax==0 ? 0 : 448/amax} 的三元式在三个 kernel 里各出现一次，
是量化代码的标准防雷点（漏掉它会产出 NaN 块，污染整个 GEMM）。

\section{前处理 Kernel 工程}\label{sec:34-kernels}

三个 kernel 全部是纯 CUDA（无 CuTe 依赖），共享 34.4~节的量化数学与一对打包
工具：

\begin{lstinputlisting}[linerange={50-82},firstnumber=auto,
  title={fp8\_gemm.cuh L50-82（量化常数、粒度枚举、Vec8 与 e4m3x2 打包）}]{../fp8_gemm.cuh}
\end{lstinputlisting}

\subsection{A 的 per-row 量化：warp-per-row}

\begin{lstinputlisting}[linerange={91-128},firstnumber=auto,
  title={fp8\_gemm.cuh L91-128（quantize\_a\_perrow\_kernel：warp 归约 + 两遍 gmem）}]{../fp8_gemm.cuh}
\end{lstinputlisting}

结构是第~\ref{ch:02}~章 warp 归约的直接应用，值得对照的是\textbf{两遍 gmem 读}：
pass 1 算行 amax，pass 2 重新读一遍行做量化写回。为什么不像第~\ref{ch:28}~章
ffpa 的 per-block kernel 那样寄存器缓存一遍读？因为那里 $K$（head\_dim）是
\textbf{编译期常量}（\texttt{kD}），整行的 Vec8 可以展开进寄存器数组；本实现的
$K$ 是运行时 GEMM 维度，寄存器数组无法定长。两遍读的代价由 L1/L2 兜底——
同一 warp 刚读过的行在 pass 2 里几乎全命中，实测该 kernel 贴内存带宽。

写侧同样关键：量化输出是 8 个 e4m3 = 8\,B，用 \texttt{uint2} 一次写——注意
这里\textbf{不能}凑 16\,B（e4m3 是 1\,B 元素，8 元素只有 8\,B；要 16\,B 得攒
16 个元素跨 lane 拼，教学版不值得）。读 16\,B 写 8\,B，字节收缩 $2\times$，
写带宽不是瓶颈。

\subsection{A 的 per-block 量化：block 二级归约}

per-block 版把统计块从一行放大到 128 行，归约升到 block 级——第~\ref{ch:03}~章
的 warp$\to$smem$\to$warp0 两级范式：

\begin{lstinputlisting}[linerange={131-170},firstnumber=auto,
  title={fp8\_gemm.cuh L131-170（quantize\_a\_perblock\_kernel：128 行块 amax 两级归约）}]{../fp8_gemm.cuh}
\end{lstinputlisting}

pass 2（L172 起，与 per-row 版逐字节同构，略）用块级 $inv\_s$ 重扫 128 行。
两个工程点：其一，块内行数=128 恰是 \texttt{kScaleBlock}：A 侧主 kernel 的
$k_{BM}{=}128$ 正好落在一个块内，B 侧 $k_{BN}{=}256$ 却横跨两块——所以
epilogue 的 scale 查表按\textbf{全局} $m/128$、$n/128$ 逐元素取
（\textbf{统计块与 tile 几何解耦}），而不是整 tile 一个标量；其二，\texttt{row >= M} 的行跳过
（尾部 block 不足 128 行），amax 不受越界垃圾影响——load 侧由 TMA 零填充兜底，
这里则必须在量化 kernel 自己 guard（它读的是原始 BF16 gmem）。

\subsection{B 的转置量化：双朝向 smem staging}\label{sec:34-bt}

三个 kernel 里最「有戏」的是 B：输出要的是 $B_8^{\top}[N,K]$ 行主序（K 内维
连续），而输入 $B[K,N]$ 的列方向 stride 是 $N$——\textbf{直接按输出朝向读写
gmem 都不 coalesced}。解法是 smem staging 双朝向 tile：

\begin{figure}[H]
\centering
\begin{tikzpicture}[
  font=\footnotesize,
  arr/.style={-{Stealth[length=2.2mm]}, thick},
  box/.style={draw, thick, minimum height=1.5cm, align=center, inner sep=4pt},
]
\node[box, draw=tkBlue, minimum width=3.0cm] (gmem) at (0,0)
  {gmem $B[K,N]$\\{\scriptsize 行主序（K 外维）}\\{\scriptsize 读：warp 沿 N 连续 16B}};
\node[box, draw=tkPurple, minimum width=3.2cm] (st) at (5.2,0)
  {smem \texttt{tile[64][128]}\\{\scriptsize 源朝向 16\,KB（bf16）}\\{\scriptsize 列方向 amax 归约}};
\node[box, draw=tkGreen, minimum width=3.2cm] (stt) at (5.2,-2.6)
  {smem \texttt{tile\_t[128][64]}\\{\scriptsize 输出朝向 8\,KB（e4m3）}\\{\scriptsize 量化散写（转置在此发生）}};
\node[box, draw=tkBlue, minimum width=3.0cm] (out) at (0,-2.6)
  {gmem $B_8^{\top}[N,K]$\\{\scriptsize 行主序（K 内维）}\\{\scriptsize 写：32 lane 各写 32B 连续}};
\draw[arr, tkBlue] (gmem.east) -- node[above, font=\tiny] {coalesced load} (st.west);
\draw[arr, tkPurple] (st.south) -- node[right, font=\tiny, align=left] {pass 2 重读\\$Q(\cdot)$ 散写} (stt.north);
\draw[arr, tkGreen] (stt.west) -- node[above, font=\tiny] {coalesced store} (out.east);
\node[font=\scriptsize, anchor=west, align=left] at (9.0,-1.3)
  {pass 1（图中上排）：逐 chunk 累积\\
   列 amax 到 \texttt{samax[128]}；\\[2pt]
   pass 2（下排）：量化 + 转置 + 写出，\\
   \texttt{samax} 原地改存 $\delta^{-1}$};
\end{tikzpicture}
\caption{B 转置量化的双朝向 staging。读路径（蓝）永远沿 B 的行（N 连续）；
写路径（绿）永远沿 $B_8^{\top}$ 的行（K 连续）；两次「转置」都发生在 smem
里（散写即转置），gmem 两侧全是 coalesced。对照第~\ref{ch:28}~章 ffpa 的
\texttt{quantize\_fp8\_vt\_kernel}：那边输出要 pad 16 列对齐 cp.async 16B，
这里 TMA 装载 OOB 自动零填、K 尾部由 store guard 裁剪，staging 免 pad。}
\label{fig:34-bt-staging}
\end{figure}

\begin{lstinputlisting}[linerange={200-240},firstnumber=auto,
  title={fp8\_gemm.cuh L200-240（quantize\_bt\_kernel：双朝向 tile 与 coalesced 装载）}]{../fp8_gemm.cuh}
\end{lstinputlisting}

装载 lambda \texttt{load\_tile} 把 B 的 $[64,128]$ 块按 16B 向量读进源朝向
tile，越界（K 尾或 N 尾）显式补零——注意这与「TMA 零填充」是两回事：这个
kernel 不走 TMA，越界补零必须自己做，但只要 amax 统计正确、写侧 guard，
零值不污染任何东西。pass 1 对每个 $n$（tid 步进 256，实际 128 个 $n$ 生效）
扫 64 行累积 \texttt{samax[n]}，跨全部 K chunk 完成整列统计。

scale 落盘段有一个值得抄走的微技巧——\textbf{\texttt{samax} 原地改存
$\delta^{-1}$}：pass 1 结束时 \texttt{samax[n]} 存的是 amax，写出
$\delta_b[n]$ 之后立刻原地覆写成 $448/\mathrm{amax}$，pass 2 查表零成本，
省一个 \texttt{inv\_s[128]} 数组与一次 \texttt{\_\_syncthreads}。per-col
关闭时（kPerCol=false），block 级再把 128 个列 amax 归约成单值写
$\delta_b[n_0/128]$，同样原地放进 \texttt{samax[0]}。

\begin{lstinputlisting}[linerange={266-288},firstnumber=auto,
  title={fp8\_gemm.cuh L266-288（pass 2：量化散写 tile\_t + 32B coalesced 写出）}]{../fp8_gemm.cuh}
\end{lstinputlisting}

pass 2 的散写循环：每个线程按 chunk 索引读 \texttt{tile[r][c8..c8+8]}，量化后
写 \texttt{tile\_t[n][r]}——下标对调即转置。写出段换成 warp 视角：一个 warp
负责一行 $n$，32 个 lane 各写 32 个连续 k 字节（\texttt{kBk t}=64 分两个
half），恰好 32B 事务；\texttt{k0+kk<K} 的 guard 处理 K 尾（该行末尾的
越界字节不写出）。两个 \texttt{\_\_syncthreads} 分别保护
tile$\to$tile\_t 的读写依赖与 tile\_t$\to$gmem 的下 chunk 覆写。

\subsection{e4m3x2 打包与饱和语义}

量化内循环的 $\hat{x}=Q(x)$ 由 \texttt{cvt\_f2\_to\_e4m3x2} 一次转换两个
float——对应 SASS/PTX 的
\texttt{cvt.rn.\allowbreak satfinite.\allowbreak e4m3x2.\allowbreak f32}
（第~\ref{ch:27}~章同指令）。\texttt{rn}
是 round-to-nearest-even，\texttt{satfinite} 把 $|x\cdot\delta^{-1}|>448$
的输入饱和到 $\pm448$。有了这条语义，量化 kernel 不需要软件
\texttt{min(max(x,-448),448)}——理论上乘 $\delta^{-1}$ 后不可能越界（$\delta$
就是按 amax 定的），但浮点除法/乘法的舍入误差可能把边界值顶出 1 ulp，
\texttt{satfinite} 恰好兜住这 1 ulp，硬件一条指令替代软件三分支。

\section{精度实测}\label{sec:34-acc}

对 $M{=}N{=}K{=}512$ 均匀随机 $[-1,1]$ 输入（srand(42)），以 CPU fp64 双精度
GEMM 为参照，四种组合 $\times$ WS/非WS $\times$ 尾部 shape 的相对 Frobenius
误差（\texttt{--fp8-gemm} 验证链的统计口径）：

\begin{center}
\footnotesize
\begin{tabular}{lcc}
\toprule
配置 & rel-Frobenius & 说明 \\
\midrule
per-row $\times$ per-col（非WS） & 0.036 & 与粒度无关的量化噪声底 \\
per-row $\times$ per-block（非WS） & 0.036 & 同上 \\
per-block $\times$ per-col（非WS） & 0.036 & 同上 \\
per-block $\times$ per-block（非WS） & 0.036 & 同上 \\
per-row $\times$ per-col（WS） & 0.036 & WS 数值等价 \\
per-block $\times$ per-block（WS） & 0.036 & WS 数值等价 \\
尾部 $M{=}300,N{=}264,K{=}144$ & 0.036 & 尾路径零额外误差 \\
\bottomrule
\end{tabular}
\end{center}

两个观察：其一，\textbf{随机数据下四种粒度重合}——均匀分布没有 outlier，
粗粒度不吃亏；真实权重/激活的分布差异才会拉开组合间的差距（这是教学实现选
随机数据的坦诚之处，也是 34.4.3~节误差模型的印证：噪声由步长 $2^{-3}$ 决定，
不由粒度决定，粒度决定的是\textbf{最坏块}的表现）。其二，\textbf{WS 与非WS
逐 bit 数值等价}——两者的 MMA 序列相同（第~\ref{ch:35}~章），差异纯粹在执行
协议。$0.036$ 落在 34.4.3~节上界 $\sqrt{2}\,\sigma_\varepsilon\approx5\%$
之内，与模型一致。

\section{小结}

本章把 GEMM 量化的「数学侧」收拢成三件事：式~\ref{eq:34-fold-rc} 的 scale
折叠恒等式（per-row/per-col 逐行列查表，per-block 每 128 行/列一个 scale）、
四种粒度的工程权衡（34.4.2~节）、以及三个前处理 kernel 的访存设计
（warp-per-row 两遍读、block 二级归约、B 的双朝向 smem 转置 staging）。
workspace 的形态已经为第~\ref{ch:35}~章铺好：$A_8[M,K]$、$B_8^{\top}[N,K]$
两个 e4m3 行主序张量 + 两个 float scale 向量，TMA descriptor 可以直接按
$(BM,BK)$/$(BN,BK)$ tile。下一章进场写主 kernel。

\section*{延伸阅读}
\addcontentsline{toc}{section}{延伸阅读}

\begin{itemize}[itemsep=1pt,topsep=2pt]
  \item 第~\ref{ch:27}~章：e4m3/e5m2 编码细节、量化算子的完整数学。
  \item 第~\ref{ch:28}~章：ffpa 的量化前处理链（per-block QK、$V^{\top}$
        转置、smoothing）——与本章对照，attention 场景的量化为什么更重
        「分布整形」。
  \item 第~\ref{ch:29}~章：$s_{dequant}=q_sk_s$ 折叠——式~\ref{eq:34-fold-rc}
        在 attention 里的带 softmax 修正版。
  \item NVIDIA TensorRT Model Quantization 的 per-channel weight
        quantization 文档：真实权重分布下 per-row/per-channel 粒度的
        校准实践。
\end{itemize}
